% =====================================================================
%  Appendix: The crest factor, the centring criterion, and diffraction
%  Self-contained. To use as an appendix inside another document,
%  delete everything above and below the marked block and \input it.
% =====================================================================
\documentclass[11pt,a4paper]{article}

\usepackage[T1]{fontenc}
\usepackage[utf8]{inputenc}
\usepackage[english]{babel}
\usepackage{lmodern}
\usepackage[a4paper,margin=2.5cm]{geometry}
\usepackage{amsmath,amssymb}
\usepackage{siunitx}
\usepackage{booktabs}
\usepackage{empheq}
\usepackage[colorlinks=true,linkcolor=black,citecolor=black,urlcolor=blue]{hyperref}

\emergencystretch=1.5em
\newcommand{\crest}{\ensuremath{C}}
\newcommand{\Uw}{\ensuremath{U_W}}
\newcommand{\Ptot}{\ensuremath{P_\mathrm{tot}}}
\newcommand{\Pavg}{\ensuremath{P_\mathrm{avg}}}
\newcommand{\Au}{\ensuremath{A_\mathrm{u}}}

\title{Appendix: The Crest Factor, the Centring Criterion,\\
       and the Diffraction Efficiency of a Multitone AOD}
\author{}
\date{}

\begin{document}
\maketitle

% ================= BEGIN APPENDIX BLOCK =================
\noindent
This appendix collects, in one place and in a single consistent notation,
the four results the main text uses: what the crest factor is
(Sec.~\ref{sec:defs}), why it acts as a hardware constraint
(Sec.~\ref{sec:constraint}), why the centring criterion forces the
\emph{worst possible} crest factor whenever an axis is driven with three
tones (Sec.~\ref{sec:centring}), and what that costs in diffracted light
(Secs.~\ref{sec:diffraction} and \ref{sec:price}).

\section{Definitions and conventions}\label{sec:defs}

\subsection{The RF signal and its envelope}

Each AOD axis is driven by one AWG channel carrying $N$ tones,
\begin{equation}
  s(t) \;=\; \sum_{n=0}^{N-1} a_n \cos\!\bigl(2\pi f_n t + \varphi_n\bigr) ,
  \qquad
  f_n \;=\; f_\mathrm{off} + n\,d ,
  \label{eq:signal}
\end{equation}
with the tone spacing $d$ and the tone phases $\varphi_n$, which are free
parameters at the AWG. Writing $\bar f$ for the mean tone frequency and
introducing the \emph{centred tone index}
\begin{equation}
  \hat n \;:=\; n-\frac{N-1}{2}
  \qquad\Longrightarrow\qquad
  f_n - \bar f = d\,\hat n ,
  \label{eq:nhat}
\end{equation}
the signal separates into a carrier at $\bar f$ and a complex envelope
\begin{empheq}[box=\fbox]{equation}
  s(t) = \operatorname{Re}\!\left[A(t)\,e^{\mathrm{i}2\pi\bar f t}\right] ,
  \qquad
  A(t) = \sum_{n=0}^{N-1} a_n\,
         e^{\,\mathrm{i}\left(2\pi d\,\hat n\,t + \varphi_n\right)} .
  \label{eq:envelope}
\end{empheq}
$A(t)$ carries everything the phases do. Its period is
$T_\mathrm{env} = 1/d$: for odd $N$ the index $\hat n$ is an integer and
$A$ itself is $T_\mathrm{env}$-periodic; for even $N$ it is half-integer and
$A(t+T_\mathrm{env}) = -A(t)$, so $A$ is \emph{anti}periodic while $|A|$ is
still periodic with $T_\mathrm{env}$. Only $|A|$ enters what follows, so the
distinction never matters -- but the averaging window below must be
$T_\mathrm{env}$ in both cases.

\subsection{Amplitude versus power: the one convention to get right}

\noindent\fbox{\parbox{0.96\textwidth}{\textbf{Two quantities in this text are
easy to confuse, and they differ by a square root.} The profile parameters
$r_x$ and $r_y$ are \emph{power} ratios -- they state how much diffracted
\emph{intensity} the edge tones receive relative to the inner ones. The
$a_n$ in Eq.~\eqref{eq:signal} are RF \emph{voltage} amplitudes. Since
diffracted intensity scales with RF power and power scales with the square
of voltage, the edge tones carry
\begin{equation}
  \varrho \;:=\; \sqrt{r_x} ,
  \qquad
  a \;=\; [\,\varrho,\;1,\;\dots,\;1,\;\varrho\,] ,
  \qquad
  P_n \;\propto\; a_n^2 .
  \label{eq:convention}
\end{equation}
At the working point $r_x = 0.97$ but $\varrho = \sqrt{0.97} = 0.984886$.
Every formula below is written in the $a_n$, i.e.\ in voltages.}}

\subsection{The crest factor}

\begin{equation}
  \crest \;:=\; \frac{\max_t |s(t)|}{\mathrm{rms}_t\,[s(t)]} .
  \label{eq:crestdef}
\end{equation}
It is a dimensionless property of the \emph{waveform}: how far the
instantaneous voltage strays above its own effective value. Both operations
run over a full envelope period $T_\mathrm{env}$.

Because the carrier at $\bar f$ oscillates much faster than the envelope,
the two time scales separate. Locally, a sinusoid of amplitude $|A|$
contributes $\tfrac12|A|^2$ to the mean square, so
$\langle s^2\rangle = \tfrac12\langle |A|^2\rangle$, while the peak of $|s|$
is the peak of $|A|$. Hence the carrier-free form
\begin{empheq}[box=\fbox]{equation}
  \crest \;=\; \sqrt{2}\;\frac{\max_t|A(t)|}{\mathrm{rms}_t|A(t)|}
         \;=\; \sqrt{2}\;\frac{\max_t|A(t)|}{\sqrt{\sum_n a_n^2}} .
  \label{eq:crestsimple}
\end{empheq}

\paragraph{The denominator does not depend on the phases.}
This is what makes the crest factor tractable. Expanding,
\begin{equation}
  \bigl\langle |A|^2 \bigr\rangle
  = \sum_{n,m} a_n a_m \Bigl\langle
      e^{\,\mathrm{i}\left[2\pi (f_n-f_m)t + \varphi_n-\varphi_m\right]}
    \Bigr\rangle_T ,
\end{equation}
and the average of a phasor over a window of length $T$ is
\begin{equation}
  \bigl\langle e^{\,\mathrm{i}[2\pi\Delta f\,t+\Delta\varphi]}\bigr\rangle_T
  = e^{\,\mathrm{i}\Delta\varphi}\,
    \frac{e^{\,\mathrm{i}2\pi\Delta f\,T}-1}{\mathrm{i}\,2\pi\Delta f\,T} ,
  \label{eq:phasoravg}
\end{equation}
which vanishes \emph{exactly} when $\Delta f\,T$ is a non-zero integer -- a
phasor completing whole turns averages to the centre of the circle
\emph{regardless of where it starts}, which is precisely why
$\Delta\varphi$ drops out. With $T = T_\mathrm{env} = 1/d$ every
$\Delta f = f_n-f_m$ is an integer multiple of $d$, so all off-diagonal
terms vanish and
\begin{equation}
  \mathrm{rms}_t|A| = \sqrt{\textstyle\sum_n a_n^2} \qquad
  \text{for \emph{every} phase set.}
  \label{eq:rmsphase}
\end{equation}
\textbf{The phases act on the crest factor through $\max_t|A|$ alone.} If the
window is chosen shorter -- a common mistake -- Eq.~\eqref{eq:phasoravg} no
longer vanishes, the double sum survives, and the rms becomes
phase-dependent.

\paragraph{Bounds.} By the triangle inequality $|A(t)|\le\sum_n a_n$ at all
times, and $|A|\ge$ its own rms somewhere, so
\begin{equation}
  \sqrt2 \;\le\; \crest \;\le\; \sqrt2\,
  \frac{\sum_n a_n}{\sqrt{\sum_n a_n^2}} \;\xrightarrow{\;a_n\equiv1\;}\;
  \sqrt{2N} .
  \label{eq:bounds}
\end{equation}
The upper bound is attained exactly when all phasors align at some instant;
the lower bound would require $|A|$ to be constant.

\section{Why the crest factor is a constraint}\label{sec:constraint}

The crest factor forbids no phase set: every $\varphi_n$ is dialable at the
AWG and the optical calculation is equally valid for all of them. What it
limits is the \emph{average RF power} the chain can still transmit linearly.

A power amplifier driven beyond its compression point clips the
\emph{instantaneous voltage}; it does not care about the time-averaged
power until thermal limits are reached, which happens much later. With
$V_\mathrm{peak} = \crest\,V_\mathrm{rms}$ and
$\Pavg = V_\mathrm{rms}^2/R$, the peak envelope power is
$P_\mathrm{peak} = \crest^2\,\Pavg$, and demanding
$P_\mathrm{peak} \le P_{1\mathrm{dB}}$ gives
\begin{empheq}[box=\fbox]{equation}
  \Pavg \;\le\; \frac{P_{1\mathrm{dB}}}{\crest^{2}} .
  \label{eq:pavglimit}
\end{empheq}
The $\SI{1}{\decibel}$ compression point, not the saturated power, is the
right threshold: linearity is exactly what the field calculation assumes.

Three budgets must be kept apart:

\begin{center}
\begin{tabular}{lll}
  \toprule
  Limit & belongs to & depends on \crest?\\
  \midrule
  \SI{2}{\watt} average RF power & AOD transducer (heating) & \textbf{no}\\
  Compression & amplifier (peak voltage) & \textbf{yes}, $\propto1/\crest^2$\\
  Diffraction saturation $\sin^2$ & AOD (total RF power) & no\\
  \bottomrule
\end{tabular}
\end{center}

\noindent With the amplifier used here, $P_{1\mathrm{dB}} = \SI{3.98}{\watt}$,
the crest factors of the working point give
$\Pavg \le \SI{0.66}{\watt}$ on the $x$ axis and \SI{0.83}{\watt} on $y$ --
well below the AOD's \SI{2}{\watt}. \emph{The amplifier is the limit, not the
AOD}, and Eq.~\eqref{eq:pavglimit} says the crest factor is the exchange rate.

\section{The centring criterion}\label{sec:centring}

\subsection{What "centred" has to mean}

The atom sits at one place, smeared by its thermal motion over
$\sigma\approx\SI{108}{\nano\meter}$. What harms it is not that the
interference pattern beats, and not that the pulse is dim: it is that the
bright spot can sit \emph{off to one side} during the pulse. The pulse area
is then larger on one edge of the position distribution than on the other,
the rotation angle depends on where the atom happened to be, and no
calibration of the mean can repair it.

The requirement is therefore not "the pattern is symmetric at some instant"
but
\begin{equation}
  I(-\mathbf{r},t) = I(\mathbf{r},t) \qquad\text{for \emph{all} } t ,
  \label{eq:centred}
\end{equation}
with the atom at $\mathbf r = 0$. This implies $\nabla I|_{0} = 0$ at all
times: the intensity has a stationary point at the atom throughout the pulse.

\subsection{The three symmetries of the profile}

The field in the focal plane is a sum over the $N_x\times N_y$ spots,
\begin{equation}
  E(\mathbf r,t) = \sum_{n,m} a_n b_m\,
    u(\mathbf r - \mathbf c_{nm})\,
    e^{-\mathrm{i}2\pi (f_n+g_m)t}\,
    e^{\,\mathrm{i}(\varphi_n+\psi_m)} ,
  \label{eq:field}
\end{equation}
with $u$ the (even) single-spot envelope, $\mathbf c_{nm}$ the spot
positions, and $g_m,\psi_m$ the $y$-axis frequencies and phases. Writing
$\bar n := N-1-n$ for the mirror index, the profile is built symmetric in
three independent ways:
\begin{align}
  \mathbf c_{\bar n\bar m} &= -\,\mathbf c_{nm}
    &&\text{(positions: the spot grid is centred),}\label{eq:symgeo}\\
  a_{\bar n} &= a_n , \quad b_{\bar m} = b_m
    &&\text{(amplitudes: } [\varrho,1,\dots,1,\varrho]\text{),}\label{eq:symamp}\\
  f_{\bar n} &= F - f_n , \quad g_{\bar m} = G - g_m
    &&\text{(frequencies).}\label{eq:symfreq}
\end{align}
Equation~\eqref{eq:symfreq} follows directly from the equidistant raster:
$f_n = f_\mathrm{off}+n\,d$ gives
$f_n + f_{\bar n} = 2f_\mathrm{off} + (N-1)d =: F$, independent of $n$.

\subsection{Derivation}

Apply the three symmetries to $E(-\mathbf r,t)$. Relabel the sum by the
mirror index, $n\to\bar n$, $m\to\bar m$ -- a permutation, so nothing is
lost:
\begin{align}
  E(-\mathbf r,t)
  &= \sum_{n,m} a_{\bar n} b_{\bar m}\,
     u(-\mathbf r - \mathbf c_{\bar n\bar m})\,
     e^{-\mathrm{i}2\pi (f_{\bar n}+g_{\bar m})t}\,
     e^{\,\mathrm{i}(\varphi_{\bar n}+\psi_{\bar m})} .
\end{align}
Now insert the three symmetries one at a time.

\paragraph{Step 1 (positions, and $u$ even).}
By Eq.~\eqref{eq:symgeo}, $-\mathbf r-\mathbf c_{\bar n\bar m}
= -\mathbf r+\mathbf c_{nm} = -(\mathbf r-\mathbf c_{nm})$, and since $u$ is
even, $u(-\mathbf r-\mathbf c_{\bar n\bar m}) = u(\mathbf r-\mathbf c_{nm})$
-- the same envelope factor as in $E(\mathbf r,t)$.

\paragraph{Step 2 (amplitudes).}
By Eq.~\eqref{eq:symamp}, $a_{\bar n}b_{\bar m} = a_n b_m$: the amplitudes
are unchanged.

\paragraph{Step 3 (frequencies).}
This is the decisive one. By Eq.~\eqref{eq:symfreq},
\begin{equation}
  e^{-\mathrm{i}2\pi (f_{\bar n}+g_{\bar m})t}
  = e^{-\mathrm{i}2\pi (F+G)t}\;
    \underbrace{e^{+\mathrm{i}2\pi (f_n+g_m)t}}_{
      \text{the \emph{conjugate} of the original factor}} .
\end{equation}
The frequency mirror does not reproduce the time factor -- it
\emph{conjugates} it, up to one global factor $e^{-\mathrm{i}2\pi(F+G)t}$
that is the same for every spot.

\paragraph{Step 4 (collect).}
Putting the three together,
\begin{equation}
  E(-\mathbf r,t) = e^{-\mathrm{i}2\pi(F+G)t}
  \sum_{n,m} a_n b_m\, u(\mathbf r-\mathbf c_{nm})\,
  e^{+\mathrm{i}2\pi(f_n+g_m)t}\,
  e^{\,\mathrm{i}(\varphi_{\bar n}+\psi_{\bar m})} ,
\end{equation}
which is to be compared with
\begin{equation}
  E^{*}(\mathbf r,t) = \sum_{n,m} a_n b_m\, u(\mathbf r-\mathbf c_{nm})\,
  e^{+\mathrm{i}2\pi(f_n+g_m)t}\,
  e^{-\mathrm{i}(\varphi_{n}+\psi_{m})} .
\end{equation}
The two sums run over the same basis functions with the same real
coefficients, and differ only in a constant phase factor per term. They are
therefore proportional -- $E(-\mathbf r,t)=e^{\mathrm{i}\alpha(t)}
E^{*}(\mathbf r,t)$, hence $I(-\mathbf r,t)=I(\mathbf r,t)$ -- if and only if
\begin{equation}
  e^{\,\mathrm{i}(\varphi_{\bar n}+\psi_{\bar m})}
  = e^{\,\mathrm{i}\Theta}\,e^{-\mathrm{i}(\varphi_{n}+\psi_{m})}
  \qquad\text{for all } n,m
\end{equation}
with a \emph{single} constant $\Theta$ independent of $n$ and $m$. Splitting
the $n$- and $m$-dependence gives the criterion:
\begin{empheq}[box=\fbox]{equation}
  \varphi_n + \varphi_{N-1-n} = c \quad\text{for all } n ,
  \qquad
  \psi_m + \psi_{M-1-m} = c' \quad\text{for all } m .
  \label{eq:centring}
\end{empheq}

\subsection{Conjugate, not mirrored}

Equation~\eqref{eq:centring} is often misread as "the phases must be
symmetric", $\varphi_{\bar n} = \varphi_n$. That is a \emph{different}
condition and it does not centre the pattern. The mirror in
Eq.~\eqref{eq:symfreq} conjugates the time factor, so the phases must
\emph{anti}-compensate, i.e.\ sum to a constant. Numerically, mirror-symmetric
phase sets leave centroid offsets of \SIrange{44}{84}{\nano\meter} against
$\sigma = \SI{108}{\nano\meter}$, while conjugate-symmetric sets give
$0.00$~nm. The two coincide only for purely real phasors,
$\varphi_n\in\{0,\pi\}$, which is the source of the confusion.

\subsection{The middle tone is its own mirror partner}\label{sec:selfpair}

A natural objection: \emph{for odd $N$ the middle tone has no partner, so its
phase is unconstrained}. It is not. The mirror acts on the indices,
$n \mapsto \bar n = N-1-n$; this is an involution, and involutions have
pairs \emph{and fixed points}:

\begin{center}
\begin{tabular}{cll}
  \toprule
  $N$ & pairing & fixed point\\
  \midrule
  3 & $0\!\leftrightarrow\!2$, \; $1\!\mapsto\!1$ & $n=1$\\
  4 & $0\!\leftrightarrow\!3$, \; $1\!\leftrightarrow\!2$ & none\\
  5 & $0\!\leftrightarrow\!4$, \; $1\!\leftrightarrow\!3$, \; $2\!\mapsto\!2$ & $n=2$\\
  \bottomrule
\end{tabular}
\end{center}

\noindent A fixed point is not an unpaired element: it is an element paired
\emph{with itself}. Equation~\eqref{eq:centring} is required for every $n$,
so at the fixed point it reads $\varphi_n+\varphi_n = c$ and
\begin{empheq}[box=\fbox]{equation}
  \varphi_{(N-1)/2} = \frac{c}{2} \quad\text{or}\quad \frac{c}{2}+\pi
  \qquad (N\ \text{odd}) .
  \label{eq:midphase}
\end{empheq}
The two roots appear because Eq.~\eqref{eq:centring} holds modulo $2\pi$.

The objection implicitly assumes each pair may carry its own constant. It may
not: $c$ \emph{is} the single global $\Theta$ of the derivation, and a single
$e^{\mathrm{i}\alpha(t)}$ has to serve all spots at once, or the
contributions do not add up to a conjugated field. The middle tone sits at
the centre frequency, so its spot lies at the mirror centre; under the mirror
it maps onto itself, and its condition is not "none" but "with itself".

\paragraph{Numerical check.} With $\varphi_x = (0,\psi,0)$ -- so $c=0$ -- and
the $y$ phases held fixed at $(0;98;16;114)^\circ$, which satisfy their own
criterion ($0+114 = 98+16$), the largest intensity gradient at the atom over
a full beat period is

\begin{center}
\begin{tabular}{crrl}
  \toprule
  $\psi$ & $\max_t|\partial_x I|/I$ & $\max_t|\partial_y I|/I$ & \\
  \midrule
  $0^\circ$   & $4\cdot10^{-10}$ & $4\cdot10^{-10}$ & criterion satisfied\\
  $10^\circ$  & $8.8\cdot10^{-4}$ & $2.6\cdot10^{-4}$ & \\
  $30^\circ$  & $2.5\cdot10^{-3}$ & $7.4\cdot10^{-4}$ & \\
  $90^\circ$  & $5.1\cdot10^{-3}$ & $1.5\cdot10^{-3}$ & \\
  $180^\circ$ & $4\cdot10^{-10}$ & $4\cdot10^{-10}$ & criterion satisfied\\
  \bottomrule
\end{tabular}
\end{center}

\noindent The gradient vanishes (to the noise floor of the numerical
derivative) at exactly the two values predicted by
Eq.~\eqref{eq:midphase} and nowhere else; with $c = 60^\circ$ the zeros move
to $\psi = 30^\circ$ and $210^\circ$. Note the third column: violating the
criterion in $x$ also tilts the $y$ gradient. With a separable (Gaussian)
envelope it would not; the real Airy envelope is radial and couples the axes,
so the criterion must hold on \emph{both}.

\subsection{Structure of the admissible family}

Writing $\varphi_n = c/2 + \chi_n$, Eq.~\eqref{eq:centring} becomes
\begin{equation}
  \chi_{\bar n} = -\,\chi_n ,
  \label{eq:antisym}
\end{equation}
i.e.\ a global phase plus an \emph{antisymmetric} part. Antisymmetric vectors
in $\mathbb{R}^N$ form a subspace of dimension $\lfloor N/2\rfloor$, so
\begin{equation}
  \dim(\text{family}) = \Bigl\lfloor \frac{N}{2} \Bigr\rfloor
  \qquad\text{instead of}\qquad N-1 .
\end{equation}
\textbf{The criterion halves the phase freedom.} One of the remaining
dimensions is moreover crest-neutral: a \emph{ramp}
$\chi_n = \delta\,\hat n$ is antisymmetric (since
$\hat{\bar n} = -\hat n$), so every ramp satisfies the criterion
automatically, and substituting it into Eq.~\eqref{eq:envelope} gives
\begin{equation}
  A_\delta(t) = A_{\delta=0}(t+\tau) , \qquad \tau = \frac{\delta}{2\pi d} ,
  \label{eq:rampshift}
\end{equation}
a pure time shift. Maximum and rms over a full period are invariant under a
shift, so a ramp cannot change \crest. Hence
\begin{equation}
  \dim(\text{crest-relevant}) = \Bigl\lfloor \frac{N}{2} \Bigr\rfloor - 1 .
  \label{eq:dimcrest}
\end{equation}
For $N=3$ this is \textbf{zero}.

\subsection{Why \boldmath$N=3$ always gives the worst crest factor}
\label{sec:n3}

The counting already says there is nothing left to vary, but the direct
calculation is shorter and needs no ramp argument at all. With
$\theta := 2\pi d\,t$ and the voltage amplitudes of
Eq.~\eqref{eq:convention}, $a_0=a_2=\varrho$, $a_1=1$:
\begin{equation}
  A(t) = \varrho\,e^{\mathrm{i}(\varphi_0-\theta)} + e^{\mathrm{i}\varphi_1}
       + \varrho\,e^{\mathrm{i}(\varphi_2+\theta)} .
\end{equation}
The criterion requires $\varphi_0+\varphi_2 = c$ and $2\varphi_1 = c$, so the
two outer phases lie symmetrically about $c/2$:
\begin{equation}
  \varphi_0 = \frac{c}{2}-\chi , \qquad
  \varphi_2 = \frac{c}{2}+\chi , \qquad
  \varphi_1 = \frac{c}{2} \;\;\text{or}\;\; \frac{c}{2}+\pi ,
\end{equation}
with $\chi$ the only remaining parameter. Factoring out $e^{\mathrm{i}c/2}$,
\begin{empheq}[box=\fbox]{align}
  A(t)\,e^{-\mathrm{i}c/2}
  &= \varrho\,e^{-\mathrm{i}(\chi+\theta)} \pm 1
   + \varrho\,e^{+\mathrm{i}(\chi+\theta)}
   \;=\; 2\varrho\cos(\theta+\chi) \;\pm\; 1 .
  \label{eq:n3env}
\end{empheq}
The sign distinguishes the two branches for $\varphi_1$. The entire time
dependence now sits in a \emph{real} cosine, which sweeps the full interval
$[-1,1]$ over one period no matter what $\chi$ is -- $\chi$ only sets
\emph{when}. Maximising over $u := \cos(\theta+\chi)\in[-1,1]$,
\begin{align}
  \text{branch }+:&\quad \max_u |\,2\varrho\,u+1\,| = 2\varrho+1
    &&\text{at } u=+1 ,\\
  \text{branch }-:&\quad \max_u |\,2\varrho\,u-1\,| = |{-}2\varrho-1| = 2\varrho+1
    &&\text{at } u=-1 ,
\end{align}
so with $a_0+a_1+a_2 = 2\varrho+1$,
\begin{equation}
  \max_t |A| = \sum_n a_n
  \qquad\text{for \emph{both} branches and \emph{every} } \chi,\,c .
\end{equation}
But $\sum_n a_n$ is the triangle-inequality bound of
Eq.~\eqref{eq:bounds}, valid for \emph{any} phase set whatsoever. The
admissible family therefore always attains it:

\medskip
\noindent\fbox{\parbox{0.96\textwidth}{\textbf{With three tones, the centring
criterion fixes the crest factor at the largest value the amplitudes allow,
for every $c$, every $\chi$ and both branches. This axis cannot be relieved
by any choice of phase.} The reason is visible in Eq.~\eqref{eq:n3env}: the
only free parameter appears \emph{inside} the cosine. It can move \emph{when}
the phasors align, not \emph{whether} they do.}}

\medskip
\noindent Numerically, with $\varrho=\sqrt{0.97}=0.9848858$:
\begin{equation}
  \sum_n a_n = 2\varrho+1 = 2.9697716 ,
  \qquad
  \sqrt{\textstyle\sum_n a_n^2} = \sqrt{2r_x+1} = \sqrt{2.94} = 1.7146428 ,
\end{equation}
-- note that the square sum contains the \emph{power} ratio linearly, the
square root of Eq.~\eqref{eq:convention} cancelling -- and therefore
\begin{equation}
  \crest_x = \sqrt2\;\frac{2.9697716}{1.7146428} = 2.4494263 .
\end{equation}
For exactly equal amplitudes this would be $\sqrt{2N}=\sqrt6 = 2.4494897$;
the working point lies $6\cdot10^{-5}$ below it, so the two agree to four
digits without being the same number.

\subsection{Contrast: \boldmath$N=4$}

Here Eq.~\eqref{eq:centring} leaves $\varphi_y = [0,\,b,\,c-b,\,c]$ with
\emph{two} parameters. The ramp is the special case $b=c/3$; any other $b$ is
an antisymmetric deviation that is not a ramp and therefore genuinely
reshapes the envelope. The two inner tones can then no longer be aligned with
the outer ones simultaneously, $\max|A|$ drops below $\sum_n a_n$, and with
it the crest factor. At $c=114^\circ$:

\begin{center}
\begin{tabular}{cccl}
  \toprule
  $b$ & $\varphi_y$ [deg] & $\max|A|$ & $\crest_y$\\
  \midrule
  $38^\circ$ & $[0,38,76,114]$ & $4.154$ & $2.8265$ \quad(ramp, $=\sum a_n$)\\
  $60^\circ$ & $[0,60,54,114]$ & $4.022$ & $2.7364$\\
  $\mathbf{98^\circ}$ & $[0,98,16,114]$ & $\mathbf{3.228}$ & $\mathbf{2.1963}$\\
  $114^\circ$ & $[0,114,0,114]$ & $3.639$ & $2.4761$\\
  \bottomrule
\end{tabular}
\end{center}

\noindent That single parameter is worth \SI{65}{\percent} more average RF
power through Eq.~\eqref{eq:pavglimit} -- and because the $x$ axis is fixed,
it is the \emph{only} adjustment the criterion leaves. The structural reason
is Eq.~\eqref{eq:dimcrest} together with the fixed-point argument: even $N$
has no self-paired tone, so no individual phase is pinned.

\section{From RF power to diffracted light}\label{sec:diffraction}

\subsection{One grating: where the \boldmath$\sin^2$ and the \boldmath$\sqrt{P}$ come from}

The RF drives a piezo transducer; the acoustic strain amplitude $S$ obeys
$P_\mathrm{ac}\propto S^2$, and the photoelastic effect turns strain into an
index modulation
\begin{equation}
  \Delta n = -\tfrac12 n_0^3 p\,S \;\propto\; \sqrt{P_\mathrm{ac}} .
  \label{eq:photoelast}
\end{equation}
\textbf{This is the origin of the square root} -- it lies in the conversion
from power to amplitude, not in the diffraction. In the Bragg regime the
incident and diffracted waves exchange amplitude periodically along the
interaction length, giving $\eta = \sin^2\nu$ with
$\nu = \pi\Delta n L/(\lambda\cos\theta_\mathrm{B})$, hence
\begin{empheq}[box=\fbox]{equation}
  \eta(P) = \sin^2\!\bigl(\kappa\sqrt{P}\bigr) ,
  \qquad
  \kappa = \frac{\pi}{\lambda}\sqrt{\frac{M_2 L}{2H}}\,\sqrt{\zeta} ,
  \label{eq:etasingle}
\end{empheq}
with the figure of merit $M_2$, the transducer aspect ratio $L/H$ and the
electro-acoustic conversion efficiency $\zeta$. None of the four is known
individually to useful accuracy, which is why $\kappa$ is \emph{measured}:
one efficiency point at a known power replaces all of them. From the test
sheet ($\SI{58}{\percent}$ in order "1-1" at \SI{1.3}{\watt} per axis, so
$\sqrt{0.58}=0.762$ per axis),
\begin{equation}
  \kappa = \frac{\arcsin\sqrt{0.762}}{\sqrt{1.3}} = \SI{0.930}{\per\sqrt\watt} .
  \label{eq:kappa}
\end{equation}

\subsection{Many tones: the coupled-mode equations}

Write the field inside the crystal as one undiffracted wave plus $N$
diffracted ones, each with a slowly varying amplitude:
\begin{equation}
  E = \Au(z)\,e^{\mathrm{i}(\mathbf k_\mathrm{u}\mathbf r-\omega_0 t)}
    + \sum_{n=0}^{N-1} A_n(z)\,e^{\mathrm{i}(\mathbf k_n\mathbf r-\omega_n t)} ,
  \qquad \omega_n = \omega_0+\Omega_n .
  \label{eq:ansatz}
\end{equation}
\emph{Notation:} \Au{} carries no number -- it belongs to no tone; $A_n$ is
the beam diffracted by tone $n$, and every $\sum_n$ below runs over the tones
and never includes \Au. Substituting into the Helmholtz equation with
$n^2 \approx n_0^2 + 2n_0\sum_n \Delta n_n\cos(\mathbf K_n\mathbf r
-\Omega_n t+\varphi_n)$, keeping terms that oscillate at $\omega_n$ and
referring them to the carrier $e^{\mathrm{i}\mathbf k_n\mathbf r}$:
\begin{align}
  \frac{\mathrm{d}A_n}{\mathrm{d}z}
    &= -\mathrm{i}\,\kappa_n e^{\mathrm{i}\varphi_n}\,\Au\,
       e^{\,\mathrm{i}\Delta k_n z} ,
  &\frac{\mathrm{d}\Au}{\mathrm{d}z}
    &= -\mathrm{i}\sum_n \kappa_n e^{-\mathrm{i}\varphi_n} A_n\,
       e^{-\mathrm{i}\Delta k_n z} ,
  \label{eq:coupledgen}
\end{align}
with $\kappa_n = \kappa\sqrt{P_n}$ and the momentum mismatch
$\Delta k_n = (\mathbf k_\mathrm{u}+\mathbf K_n-\mathbf k_n)\cdot\hat{\mathbf z}$.
Three things follow immediately.

\paragraph{(a) The tone phases drop out.} The substitution
$\tilde A_n = A_n e^{-\mathrm{i}\varphi_n}$ removes $\varphi_n$ from both
equations, and $|A_n| = |\tilde A_n|$. \textbf{The diffraction efficiency is
exactly independent of the tone phases}; $\varphi_n$ reappears only as the
phase of the diffracted wave, i.e.\ in the interference pattern, where it
belongs. This is why the crest factor -- a property of the \emph{sum} $s(t)$,
which only the broadband amplifier ever sees -- cannot appear anywhere in
this section.

\paragraph{(b) Phase matching is assumed here, not later.} Setting
$\Delta k_n = 0$ and transforming the phases away leaves
\begin{equation}
  \frac{\mathrm{d}\Au}{\mathrm{d}z} = -\mathrm{i}\sum_n \kappa_n A_n ,
  \qquad
  \frac{\mathrm{d}A_n}{\mathrm{d}z} = -\mathrm{i}\,\kappa_n \Au ,
  \label{eq:coupled}
\end{equation}
where $z$ is measured in units of the interaction length, so the crystal exit
is at $z=1$ and $\kappa_n$ is the coupling accumulated over the whole
crystal. At the working point the tone span is \SI{0.37}{\mega\hertz} against
a \SI{36}{\mega\hertz} bandwidth, worth \SI{0.007}{\percent} in efficiency;
an \SI{18}{\mega\hertz} span would already cost \SI{15}{\percent}.

\paragraph{(c) Unequal $\Delta k_n$ would break the reduction below}, because
carrying $e^{-\mathrm{i}\Delta k_n z}$ through the derivative produces a new
combination with weights $\kappa_n\Delta k_n$, and the hierarchy does not
close unless all $\Delta k_n$ are equal (or zero).

\subsection{Reduction to two waves}

\paragraph{Bright mode.} Define
\begin{equation}
  \kappa_\mathrm{eff} = \sqrt{\textstyle\sum_n \kappa_n^2} ,
  \qquad
  B = \frac{1}{\kappa_\mathrm{eff}}\sum_n \kappa_n A_n .
  \label{eq:brightmode}
\end{equation}
Then $\mathrm{d}\Au/\mathrm{d}z = -\mathrm{i}\kappa_\mathrm{eff}B$ and
$\mathrm{d}B/\mathrm{d}z = -\mathrm{i}\kappa_\mathrm{eff}\Au$: a closed
two-wave system.

\paragraph{Dark modes stay dark.} Any $C=\sum_n c_n A_n$ with
$\sum_n c_n\kappa_n = 0$ obeys $\mathrm{d}C/\mathrm{d}z = 0$, and $C(0)=0$
since nothing is diffracted at the entrance. So $C\equiv0$: of the $N$
diffracted degrees of freedom, $N-1$ are never excited. This is what makes
the reduction exact rather than approximate.

\paragraph{Solve and return.} With $\Au(0)=1$, $B(0)=0$:
$\Au = \cos(\kappa_\mathrm{eff}z)$, $B = -\mathrm{i}\sin(\kappa_\mathrm{eff}z)$.
Integrating the second of Eq.~\eqref{eq:coupled} with $A_n(0)=0$, the same
integral appears for every $n$:
\begin{equation}
  A_n(z) = -\mathrm{i}\,\kappa_n\!\int_0^z\!\Au(z')\,\mathrm{d}z'
         = -\mathrm{i}\,\frac{\kappa_n}{\kappa_\mathrm{eff}}
           \sin(\kappa_\mathrm{eff}z) ,
\end{equation}
so the orders differ only by the prefactor $\kappa_n$. Summing at $z=1$,
\begin{equation}
  \eta_\mathrm{tot} = \sum_n |A_n|^2 = \sin^2\kappa_\mathrm{eff}
  \;=\; 1-|\Au|^2 ,
\end{equation}
and -- the step where the sum moves under the root --
\begin{equation}
  \kappa_\mathrm{eff} = \sqrt{\textstyle\sum_n \kappa_n^2}
                      = \sqrt{\textstyle\sum_n \kappa^2 P_n}
                      = \kappa\sqrt{\Ptot} .
  \label{eq:keffP}
\end{equation}
The squares of the couplings \emph{are} the powers, so their quadratic sum is
nothing but the ordinary sum of powers under one square root. Hence
\begin{empheq}[box=\fbox]{equation}
  \eta_\mathrm{tot} = \sin^2\!\bigl(\kappa\sqrt{\Ptot}\bigr) ,
  \qquad
  \eta_n = \frac{P_n}{\Ptot}\,\eta_\mathrm{tot} ,
  \qquad \Ptot = \sum_n P_n .
  \label{eq:multitone}
\end{empheq}
Energy is conserved exactly, $|\Au|^2+\sum_n|A_n|^2 = 1$, and for $N=1$
Eq.~\eqref{eq:multitone} reduces to Eq.~\eqref{eq:etasingle}.

\paragraph{The sum belongs under the root.} The $N$ gratings act as
\emph{one} grating of strength $\sqrt{\sum_n\kappa_n^2}$: an additional tone
does not open a new channel out of the undiffracted beam, it shares the
existing one. The naive alternative $\sum_n\sin^2(\kappa\sqrt{P_n})$ is not
wrong physics but the small-signal limit taken in the wrong place -- for
$\kappa\sqrt{P}\ll1$ both expressions tend to $\kappa^2\Ptot$. They separate
exactly where we work:

\begin{center}
\begin{tabular}{cccc}
  \toprule
  \Ptot & $\kappa\sqrt{\Ptot}$ & $\eta_\mathrm{tot}$ exact & naive sum (3 tones)\\
  \midrule
  \SI{5}{\milli\watt}   & $0.066$ & $0.00432$ & $0.00433$ \;($+0.10\,\%$)\\
  \SI{50}{\milli\watt}  & $0.208$ & $0.0427$  & $0.0431$ \;($+0.97\,\%$)\\
  \textbf{\SI{0.664}{\watt}} & $0.758$ & $\mathbf{0.473}$ & $0.539$ \;($\mathbf{+14.0\,\%}$)\\
  \bottomrule
\end{tabular}
\end{center}

\noindent A useful consequence: \emph{the total diffraction efficiency
depends only on the total RF power, not on how it is divided among tones.}
Splitting into a multitone profile costs no efficiency at all.

\section{What the crest factor costs in diffracted light}\label{sec:price}

The two halves now join. Equation~\eqref{eq:multitone} says the diffraction
depends only on \Ptot; Eq.~\eqref{eq:pavglimit} says the crest factor caps
\Ptot. Substituting one into the other gives the whole optical effect of the
crest factor in a single expression:
\begin{empheq}[box=\fbox]{equation}
  \eta_\mathrm{axis}(\crest)
  = \sin^2\!\left(\frac{\kappa\sqrt{P_{1\mathrm{dB}}}}{\crest}\right) ,
  \qquad
  \eta_\mathrm{profile} = \eta_x\,\eta_y .
  \label{eq:etacrest}
\end{empheq}
With $\kappa = \SI{0.930}{\per\sqrt\watt}$ and
$P_{1\mathrm{dB}} = \SI{3.98}{\watt}$ the numerator is
$\kappa\sqrt{P_{1\mathrm{dB}}} = 1.8559$, so $\eta_\mathrm{axis} =
\sin^2(1.8559/\crest)$:

\begin{table}[htb]
  \centering
  \caption{The crest factor translated into diffraction, via
  Eq.~\eqref{eq:etacrest}.}
  \begin{tabular}{lcccl}
    \toprule
    & \crest & $\Pavg = P_{1\mathrm{dB}}/\crest^2$ & $\eta_\mathrm{axis}$ & \\
    \midrule
    theoretical minimum & $\sqrt2 = 1.414$ & \SI{1.99}{\watt} & $0.935$ & unreachable\\
    $x$ with $\psi=90^\circ$ & $1.822$ & \SI{1.20}{\watt} & $0.725$ & centring abandoned\\
    $y$, $b = 98^\circ$ & $2.196$ & \SI{0.83}{\watt} & $0.559$ & working point\\
    $\mathbf{x}$, \textbf{centred,} $\mathbf{N=3}$ & $\mathbf{2.449}$ & \textbf{\SI{0.66}{\watt}} & $\mathbf{0.472}$ & \textbf{forced, Sec.~\ref{sec:n3}}\\
    $y$, ramp & $2.827$ & \SI{0.50}{\watt} & $0.373$ & worst case\\
    \bottomrule
  \end{tabular}
\end{table}

\noindent At the working point,
$\eta_\mathrm{profile} = 0.472\times0.559 = \SI{26.4}{\percent}$, i.e.\
between \SI{2.0}{\percent} and \SI{2.4}{\percent} per spot over the twelve
spots. Had both axes reached the theoretical minimum $\crest=\sqrt2$, the
profile efficiency would be \SI{87}{\percent} -- a factor $3.3$ more light.
That factor is the price of the crest factor, and Sec.~\ref{sec:n3} says how
much of it is negotiable: on the $x$ axis, none of it, as long as three tones
are used and the pattern must stay centred on the atom.

\medskip
\noindent Three closing remarks.
\begin{enumerate}\itemsep0.2em
  \item \textbf{The chain is one-directional.} Phases $\to$ \crest{} $\to$
        \Pavg{} $\to$ $\eta$ $\to$ Rabi frequency. Nowhere do the phases act
        on $\eta$ directly; see Sec.~\ref{sec:diffraction}(a). The crest
        factor is a budget limit of the electronics that happens to
        propagate all the way to the atom.
  \item \textbf{The remedy is more tones, not better phases.} By
        Eq.~\eqref{eq:dimcrest} the crest-relevant freedom is
        $\lfloor N/2\rfloor-1$, zero only for $N\le3$. A fourth $x$ tone
        dissolves the conflict instead of negotiating it.
  \item \textbf{One channel is not covered here.} Equation~\eqref{eq:etasingle}
        assumes a linear transducer. The instantaneous strain in the crystal
        carries the crest factor \crest{} as well, so a high \crest{} also
        brings on acoustic nonlinearity earlier -- an effect independent of
        the amplifier, pointing the same way, and not quantified here.
\end{enumerate}
% ================== END APPENDIX BLOCK ==================

\end{document}
